% Gráfico 8. Tipo de cambio efectivo real y exportaciones en volumen (2007 = 100), 2007-2025.
% Datos: ../datos/g08-competitividad.csv (Eurostat ert_eff_ic_a y nama_10_gdp). Un aumento del tipo de cambio efectivo real es una apreciación.
\pgfplotstableread[col sep=comma,comment chars={\#}]{datos/g08-competitividad.csv}\gOcho
\begin{tikzpicture}[tcee]
  \begin{axis}[dicgraf,
      xmin=2006.6, xmax=2025.4, xtick={2007,2009,...,2025},
      ymin=80, ymax=160, ytick={80,100,...,160},
      legend columns=2, dicapas]
    \capa{1}{\dicraya[dicumbral]{(2006.6,100) (2025.4,100)}}
    \capa{2}{\addplot[diclinea=dicespana, leyendalinea] table[x=anio, y=tcer_ipc42] {\gOcho}; \addlegendentry{Tipo de cambio efectivo real (IPC, 42 socios)}}
    \capa{3}{\addplot[diclinea=dicespana, densely dashed, leyendalinea] table[x=anio, y=tcer_clu_ze] {\gOcho}; \addlegendentry{Tipo de cambio efectivo real (CLU, zona del euro)}}
    \capa{4}{\addplot[diclinea=dicexterna, leyendalinea] table[x=anio, y=export_real] {\gOcho}; \addlegendentry{Exportaciones de bienes y servicios en volumen}}
  \end{axis}
\end{tikzpicture}
